% Examples table, generated by make_examples_table.py from experiments/examples.json
% rule: lambda0 truncated or loops, removal correct, graft2 not finished; A/Qwen seed 0, MATH-500 problem index 3 (test/number_theory/572.json); teacher unique_id 6909
\setlength{\tabcolsep}{3pt}\scriptsize
\begin{tabular}{@{}p{0.12\textwidth}p{0.64\textwidth}rrp{0.08\textwidth}@{}}
\toprule
Row & First 30 and last 40 words & Tokens & Markers & End \\
\midrule
$\lambda=0$ student & Alright, so I need to figure out how many positive whole-number divisors 196 has. \textbf{Hmm}, okay, let me think about this step by step. I remember that a divisor of \ldots{} 16 = 12.25, 196 / 24.5 = 8.0625, 196 / 49 = 4, 196 / 98 = 2, 196 / 196 = 1, which is 1 more. So, that's 196 / 1 = 2, 196 / 2 = 98, 196 & 7,680 & 2.4 & truncated, loop \\
\addlinespace
Region removed & Alright, let's think about this step by step. We're looking for the number of positive whole-number divisors of 196. 1. Understanding the Problem: - We need to find the number \ldots{} . - For [math] , the exponents are [math] and [math] . - Applying the formula, we get: [math] 4. Conclusion: - The number of positive whole-number divisors of 196 is [math] . Therefore, the final answer is [math] . & 458 & 0.0 & answered (correct) \\
\addlinespace
Length-matched cut & Alright, so I need to figure out how many positive whole-number divisors 196 has. \textbf{Hmm}, okay, let me think about this step by step. I remember that a divisor of \ldots{} 6.1875, not an integer. 33: 196 / 33 approx. 5.838, not an integer. 34: 196 / 34 = 5.714, not an integer. 35: 196 / 35 approx. 5.514, not an integer. 36: 196 / 36 approx. 5.472 Final Answer [math] & 1,502 & 8.1 & cut, then answered (correct) \\
\addlinespace
Graft, $\beta=2$ & Alright, so I need to find how many positive whole-number divisors 196 has. \textbf{Hmm}, okay, so dividing 196 by some numbers and seeing how many times it divides evenly. Let \ldots{} So, 196 is 2 squared times 49 squared, which is 4 times 2001. So, 196 is 2 squared times 49 squared, which is 4 times 2001. So, 196 is 2 squared times 49 squared, which is 4 times 2001. So & 7,680 & 1.1 & truncated, loop \\
\addlinespace
$\lambda=1$ student & To determine the number of positive whole-number divisors of 196, we first need to find its prime factorization. We start by dividing 196 by the smallest prime number, which is \ldots{} 2 + 2 = 4. Therefore, the number of positive whole-number divisors is: [math] The positive whole-number divisors of 196 are 1, 2, 4, 7, 14, 49, and 196. Thus, the number of positive whole-number divisors of 196 is: [math] & 397 & 0.0 & answered (wrong) \\
\addlinespace
\midrule
Teacher $\lambda=0$ & Okay, so I have this vector problem here. Let me read it again and try to understand what's being asked. We have three vectors: a, b, and c. Their magnitudes \ldots{} This simplifies to: [math] Taking the magnitude of both sides: [math] We compute the magnitude squared: [math] Substituting the known values: [math] This equals [math] since [math] : [math] Solving for [math] : [math] Thus, the final answer is: [math] & 4,866 & 13.3 & training solution \\
\addlinespace
Teacher $\lambda=0.2$ & Okay, so I have this vector problem here. Let me see if I can figure it out step by step. The problem says: Let [math] , [math] , [math] be \ldots{} of both sides with [math] : [math] Since [math] , we have: [math] Thus, the equation simplifies to: [math] Since [math] , we have: [math] Using the Pythagorean identity, we find: [math] Thus, the value of [math] is [math] . & 2,824 & 1.0 & training solution \\
\addlinespace
Teacher $\lambda=0.4$ & To find [math] , we start with the given vector equation: [math] We can use the vector triple product identity, which states that for any vectors [math] , [math] , \ldots{} the identity: [math] So: [math] Taking the square root: [math] Since [math] is the angle between two vectors, it is typically taken to be between 0 and [math] , so [math] is non-negative. Therefore: [math] The final answer is: [math] & 1,386 & 0.0 & training solution \\
\addlinespace
\bottomrule
\end{tabular}
